\section{Problem Setting}
\label{sec:problem-setting}

\paragraph{Classifiers, scores and shift.}
A fine-tuned classifier maps an input $x$ to logits $z(x) \in \mathbb{R}^K$ over $K$ classes ($K = 2$ for offensive
language, $K = 3$ for sentiment). It is trained on a labelled source distribution $P$ and deployed on a target
distribution $Q$ of a different language variety: a regional dialect, a romanised or code-mixed register, or another
language. We observe the source training, validation and calibration splits and, on the target, an unlabeled stream
and possibly a small number $n$ of labelled calibration items. All target quantities are evaluated on a disjoint target
test split.

\paragraph{Post-hoc corrections.}
Every correction we study is a map applied to the logits of a fixed model. Temperature scaling
(TS; \citealp{guo2017calibration}) uses $\mathrm{softmax}(z/T)$ with $T$ fitted on the source calibration split
(\emph{source TS}) or on target labels. TS is monotone in each logit difference and cannot change the predicted class.
Bias-corrected temperature scaling (BCTS; \citealp{alexandari2020maximum}) uses $\mathrm{softmax}(z/T + b)$ and can
move the decision threshold. With $n$ target labels we fit both by maximum a posteriori estimation around the source
estimate (\emph{shrinkage TS} and \emph{MAP BCTS}; Gaussian priors whose variances are set from oracle deviations
measured on the other task, so nothing is fitted on the evaluated target). Batch calibration
\citep{zhou2024batch} needs no labels: it divides the source-TS probabilities by their mean over an unlabeled target
batch and renormalises, which equalises the predicted class marginals and therefore assumes a uniform target class
prior. Its prior-aware variant multiplies by a class prior estimated from the $n$ labels.

\paragraph{Quantities that describe a shift.}
For a model and a target we record, after source TS:
(i) the \emph{ranking} AUROC of the scores (macro one-vs-rest for $K = 3$) and its \emph{retention}, the ratio to the
in-domain AUROC;
(ii) the \emph{oracle temperature ratio} $\log(T^{*}_{Q} / T_{P})$, where $T^{*}_{Q}$ maximises the likelihood of all
target calibration labels (we flag fits that reach the upper bound $T = 20$, i.e.\ an unbounded optimum);
(iii) the \emph{explained label entropy} $1 - \mathrm{NLL}_{Q}(\text{target TS}) / H(\pi_Q)$, where $\pi_Q$ is the
target class prior; values $\le 0$ mean that even the oracle temperature is no better than predicting $\pi_Q$ for every
input;
(iv) the \emph{predicted-marginal shift}, the total variation between the predicted-label distribution on the target
inputs and on the source calibration split, which needs no target labels.

\paragraph{Regimes.}
These quantities separate four regimes, which we named before the cross-source experiment:
\emph{no correction needed} (ranking and calibration transfer);
\emph{temperature shift} (a finite change of the optimal temperature with a stable predicted marginal);
\emph{threshold shift with retained ranking} (the ranking survives, but the predicted marginal moves and the decision
threshold is wrong); and
\emph{discrimination collapse} (the score carries no usable information about the target label; the oracle temperature
diverges and target TS is no better than a constant). We add a fifth, \emph{weak ranking}, for targets where the ranking
is above chance but target TS is still no better than a constant. The operational thresholds that assign targets to
regimes were set after inspecting the results and are reported as such (Section~\ref{sec:results}).

\paragraph{Why temperature can be the wrong knob.}
Under collapse no post-hoc map of the score can beat the constant predictor at the target prior
(Proposition~\ref{prop:impossibility}, Section~\ref{sec:shift-triage}), and abstaining on low-confidence inputs leaves
the selective risk at the majority-class error. Under a threshold shift the ranking is intact but every
temperature-only correction leaves the decisions unchanged; only a bias term or a different system can repair them.
Distinguishing these regimes requires labels for collapse but not for presence of shift, which motivates the
label-efficient triage procedure of Section~\ref{sec:shift-triage}.
